problem:  
Let \(M^n\) be a closed, simply connected Riemannian manifold with **nonnegative sectional curvature**, admitting an **isometric fixed point homogeneous \(S^1\)-action**.  
Let \(F\subset \mathrm{Fix}(S^1,M)\) denote the maximal-dimensional fixed point component, and let \(Q=M/S^1\) be the orbit space endowed with its induced Alexandrov metric of curvature \(\ge 0\). Then \(\partial Q=F\), and the distance-to-boundary function  
\[
r(x)=\operatorname{dist}(x,F)
\]
is concave on \(Q\). Denote by \(E\subset Q\) the **extremal set** consisting of points where \(r\) attains its maximal value.  

**Problem:**  
Does \(Q\) necessarily **admit an Alexandrov join decomposition**  
\[
(Q,d_Q) \cong (F,d_F) * (E,d_E)
\]
for some intrinsic metrics \(d_F\) on \(F\) and \(d_E\) on \(E\) (both Alexandrov spaces with curvature \(\ge 0\)), satisfying the expected dimension condition  
\(\dim(F)+\dim(E)+1 = \dim(Q)\)?  
Equivalently, is every such orbit space \(Q=M/S^1\) with nonnegative curvature **globally an Alexandrov join** between its boundary \(F\) and its maximal-distance extremal stratum \(E\)?  

If not, can one construct or exclude an example where \(M\) is smooth and nonnegatively curved, but \(Q\) fails to admit such a decomposition even topologically?  

Why is it a "good" problem:  
This formulation isolates a **precisely defined metric–topological rigidity question** at the junction of **Riemannian transformation group theory** and **Alexandrov geometry**.  

- **Foundational depth:**  
  For fixed point homogeneous manifolds, known results (Grove–Searle, Grove–Wilking) guarantee a double disk-bundle decomposition of \(M\), but not a global metric characterization of \(Q\). The question asks whether the combination of **nonnegative curvature** and **fixed point homogeneity** automatically imposes an **Alexandrov join structure** on the orbit space \(Q\)—a metric analogue of the double disk-bundle topology.  

- **Rigidity vs. Flexibility frontier:**  
  If true, it would yield a new **global metric rigidity phenomenon**, asserting that symmetry and curvature together fix the entire metric geometry of the quotient. If false, the construction of a counterexample would reveal new **flexibility mechanisms** for nonnegatively curved manifolds with large symmetry, refining our understanding of curvature–symmetry interactions.  

- **Cross-field relevance:**  
  The problem interlaces tools from:
  - **Alexandrov geometry** (concave distance functions, extremal subsets, join structures);  
  - **Riemannian geometry and group actions** (fixed point homogeneity, double disk-bundle decompositions);  
  - and potentially **geometric analysis** (via generalized Sharafutdinov retractions).  
  A resolution would deepen understanding of how *smooth curvature constraints* manifest as *metric singular structures* in quotients.  

Thus it is a *good* problem: sharply framed, conceptually transparent, and poised at a major open interface—whether **global Alexandrov join rigidity** is an intrinsic feature of nonnegatively curved, fixed point homogeneous manifolds.